\chapter{Results}\label{ch:results}

This chapter reports what was measured, in mechanism order: first what
the judge does, then what the partition does.

That ordering is load-bearing rather than editorial. The strongest
result in this project is a measurement of the judge, and it is what
makes the partition results of Section~\ref{sec:res-partition}
interpretable instead of merely disappointing. A reader who learns in
Section~\ref{sec:res-judge} that the judge's verdict tracks answer
correctness already knows why the rubric contrast in
Section~\ref{sec:res-partition} is flat, and knows that the flatness is
a property of the corpus and the judging paradigm rather than a failure
of the construction.

The reading conventions below govern every result in this chapter.
Each is stated once here and referenced thereafter; none is weakened
by not being repeated at the point of use.

\begin{center}
\fbox{\parbox{0.92\textwidth}{%
\textbf{Reading conventions for every result in this chapter.}
\begin{itemize}
\item \textbf{Tie policy.} Every Spearman $\rho$ appears under both
tie conventions --- ties half-weighted, then tied comparisons dropped,
in that order --- per rule~D5 of Section~\ref{sec:meth-discipline};
the measured instability that forces the rule is
Section~\ref{sec:res-tie-policy}. Every table keeps both rows.
\item \textbf{Estimator.} Bradley--Terry $\theta$ is quoted, never
win-rate, because the scheduler over-samples uncertain pairs and
win-rate is biased by opponent strength
(Section~\ref{sec:meth-scheduling}).
\item \textbf{Seed.} All construction and arena results are at seed
42. There is no three-seed convention behind any headline number in
this thesis, and none is claimed.
\item \textbf{Counts.} Comparison counts are logical,
$B_{\text{logical}}$; each logical comparison issues two judge calls
under the dual-call design, so $B_{\text{API}} = 2B_{\text{logical}}$.
\item \textbf{Quantisation.} Spearman $\rho$ over $M$ models moves on
a grid of $6/(M(M^2-1))$: $0.0119$ at $M = 8$, $0.0046$ at $M = 11$,
$0.0035$ at $M = 12$. Differences below one grid step are not
differences.
\item \textbf{Evidence marks.} Every substantive claim is marked
\textsc{measured} or \textsc{argued} in the sense fixed in
Section~\ref{sec:research-questions}.
\end{itemize}}}
\end{center}

Interpretation is deferred to Chapter~\ref{ch:discussion}.

\section{What Was Run}\label{sec:res-runs}

Three arena runs and five offline analyses produced the results in this
chapter.

\textbf{Run~A --- the pilot Math arena.}
Snapshot \texttt{20260727\_042523} (154 nodes, 87 leaves, certified;
87/87 leaves carrying rubrics), 8 models, 11 Math cells, 241 judged
held-out questions, 1{,}736 verdicts per condition, two conditions ---
\texttt{MAIN} (cell-scoped reference-informed rubric) and
\texttt{GENERIC\_JUDGE} (same cells, domain-agnostic prompt). Completed
with no errors beyond a known inverse-variance floor guard. This run
used the centrality-ordered query sampler and predates the corpus
repair described in Section~\ref{sec:res-capture-gap}; it carries the
capture-gap defect reported there.

\textbf{Run~B --- the paired no-partition baseline on Math.}
The 12-model length-matched band on the repaired corpus, judging an
\emph{identical} set of 250 held-out Math questions in both arms:
\texttt{MAIN} (each question judged under its own leaf's rubric,
Bradley--Terry fitted per leaf and rolled up by pooling) against
\texttt{C5} (the same questions judged under a generic MT-Bench-style
pairwise prompt with no partition at all, Bradley--Terry fitted
directly at the domain). 5{,}278 comparisons. Shuffled query sampling.
This is the no-taxonomy baseline, and it is the run that tests the
partition itself rather than the rubric.

\textbf{Run~C --- law, on the Run~B roster.} Both arms, the domain the
offline screen selects (Section~\ref{sec:res-law}); same twelve models,
same paired design. Two scheduler changes distinguish it from Run~B and
are part of the record: a mandatory bootstrap phase gave every one of
the 66 pairs at least two comparisons before any adaptive scheduling
(Run~B left 46 of 66 pairs with no data), and per-(leaf, pair) query
shuffling raised held-out pool coverage to 246 of 297 leaf-routed law
questions (83\%, against 66\% in both Math runs). 4{,}978 comparisons.
Results in Section~\ref{sec:res-c5}.

\textbf{The offline analyses} required no judge calls: the fixed-point
certificate and the acceptance-gate seed sweep
(Section~\ref{sec:res-construction}); the rubric-specificity null
(Section~\ref{sec:res-rubric-null}); the pre-registered granularity
analysis (Section~\ref{sec:res-granularity}); the domain-reordering
screen (Section~\ref{sec:res-law}); and the reliability-constant refit
reported with the measurement debt in Section~\ref{sec:disc-debt}.

\section{What the Judge Does}\label{sec:res-judge}

\subsection{The Judge Decides on Answer Correctness}\label{sec:res-correctness}

\textsc{measured.} On Run~A, restricted to items where exactly one of
the two responses is correct and the judge did not return a tie, the
judge selects the correct response $914/1029 = 88.8\%$ of the time
under \texttt{MAIN} and $913/1033 = 88.4\%$ under
\texttt{GENERIC\_JUDGE}. The leaf-specific rubric moves that number by
$0.4$ percentage points.

\textsc{measured}, and not part of the hypothesis: the tie rate roughly
triples exactly where the answer key stops discriminating.

\begin{table}[h]
\centering
\caption{Tie rate by whether the answer key can discriminate between
the two responses (Run~A, both conditions). The subsets partition the
1{,}736 comparisons of each condition.}
\label{tab:tie-by-key}
\begin{tabular}{lrrr}
\toprule
Subset & $n$ & tie rate \texttt{MAIN} & tie rate \texttt{GENERIC} \\
\midrule
exactly one response correct & 1160 & $\sim$11\% & $\sim$11\% \\
both responses correct       & 129  & \textbf{38.0\%} & \textbf{32.6\%} \\
neither response correct     & 447  & \textbf{33.1\%} & \textbf{34.9\%} \\
\bottomrule
\end{tabular}
\end{table}

A three-fold jump in ties precisely where the key stops helping is the
signature of a correctness-driven decision. It arrived from a direction
nobody was looking in: the tie rate was recorded for the
Bradley--Terry sufficient statistics, not as a test of anything.

Two qualifications travel with the $88.8\%$ figure and are stated here
rather than later. First, half of its support is format-decided
(Section~\ref{sec:res-capture-gap}). Second, the judge is
answer-key-blind but not answer-blind: it sees the multiple-choice
options, and $94.8\%$ of stored traces state the answering model's own
selection (Section~\ref{sec:arch-judge}). Nothing in the prompt tells
the judge which option is right; a judge able to solve the item does
not need to be told.

\paragraph{Fidelity when the key cannot discriminate.}
The natural follow-up is to recompute rank fidelity on the 576
comparisons (170 questions) where correctness cannot discriminate ---
both responses correct, or neither --- so that the judge must evaluate
rather than verify. Two components of that analysis are secure and one
is not.

Secure, and \textsc{measured}: the tie fraction rises from $0.11$ on
the discriminable subset to $0.34$ on the correctness-blind one, in
both conditions. And the \textbf{selection control holds}: on those same
576 rows the ground-truth model ordering is \emph{identical} to the
full-set ordering, $\rho(\text{GT}_{\text{blind}},
\text{GT}_{\text{full}}) = +1.000$, with ground-truth accuracy spread
$0.249$ against $0.258$ on the full set --- a $3.5\%$ reduction. The
correctness-blind subset is therefore not range-restricted, and any
fidelity drop on it is judgment rather than selection. That control was
fixed in advance and it is the part of this analysis that survives
unambiguously, because it touches ground-truth accuracy only and never
a verdict.

Not secure: the magnitude of the fidelity drop itself. The value
originally recorded --- Spearman $\rho$ falling from the $0.95$--$0.98$
region to $0.74$--$0.76$ --- \textbf{could not be reproduced}:
recomputation on the same verdict set returns $0.90$--$0.93$ under
both tie conventions, and the two computations differ by an analysis
convention that was not written down at the time. The magnitude is
therefore withdrawn, and what carries the claim is the direction and
the tie behaviour: fidelity on the correctness-blind subset is lower
than on the full set, and the judge abstains on roughly a third of
those comparisons. The reconstruction detail and its resolution path
are recorded in the corrections register
(Appendix~\ref{app:corrections-register}).

\subsection{Reasoning Text Makes the Judge Worse at Matching the Key}\label{sec:res-trace-inversion}

\textsc{measured.} Stratifying Run~A's decidable, non-tied comparisons
by whether each response carries a chain-of-thought trace inverts the
design assumption. Agreement with the answer key is \emph{highest}
where neither response carries any reasoning and \emph{lowest} where
both do.

\begin{table}[h]
\centering
\caption{Answer-key agreement by trace stratum (Run~A). Rows where
exactly one response is correct and the judge did not tie. The roster
is balanced 4/4 on trace presence.}
\label{tab:trace-strata}
\begin{tabular}{llrrrr}
\toprule
Condition & Stratum & hits\,/\,$n$ & agreement & SE & ties \\
\midrule
\texttt{MAIN} & trace $\times$ trace & 181\,/\,215 & \textbf{84.2\%} & 2.5\% & 40 \\
\texttt{MAIN} & trace $\times$ no-trace & 590\,/\,660 & 89.4\% & 1.2\% & 17 \\
\texttt{MAIN} & no-trace $\times$ no-trace & 143\,/\,154 & \textbf{92.9\%} & 2.1\% & 74 \\
\addlinespace
\texttt{GENERIC} & trace $\times$ trace & 174\,/\,211 & \textbf{82.5\%} & 2.6\% & 44 \\
\texttt{GENERIC} & trace $\times$ no-trace & 595\,/\,670 & 88.8\% & 1.2\% & 7 \\
\texttt{GENERIC} & no-trace $\times$ no-trace & 144\,/\,152 & \textbf{94.7\%} & 1.8\% & 76 \\
\bottomrule
\end{tabular}
\end{table}

The gap is $8.7$ percentage points under \texttt{MAIN}, about $2.7$
standard errors, with the same shape under \texttt{GENERIC}.
The tie counts corroborate it independently and by a different
mechanism: 74 ties in the no-trace stratum against 17 in the mixed one.
Given two bare answer letters and no reasoning to compare, the judge
either matches the key or declares equivalence.

This is only paradoxical if the judge is evaluating. It is exactly what
verification predicts, with the prose acting as a distraction that
occasionally argues the judge out of the correct answer. Like the tie
result above, it was not part of the hypothesis, which is why it counts
as corroboration rather than as fit.

\begin{figure}[h]
\centering
\fbox{\parbox{0.9\textwidth}{\centering\vspace{0.6cm}
Figure --- Answer-key agreement by trace stratum (placeholder)\\
\footnotesize Grouped dot-with-CI plot: three trace strata on the
horizontal axis, agreement with binomial CI on the vertical, one series
per condition; a second panel repeats the display at matched capability
gap (Table~\ref{tab:capture-survives}); data: Run~A verdict export
joined to the response-format classification.\vspace{0.6cm}}}
\caption{Answer-key agreement by trace stratum, with the capability-gap
control alongside. The two panels are the same measurement before and
after conditioning on the ground-truth accuracy gap between the two
models being compared.}
\label{fig:trace-strata}
\end{figure}

\subsection{A Defect Found by an Independent Route: The Capture Gap}\label{sec:res-capture-gap}

\textsc{measured.} Four of Run~A's eight models had their stored
response text recorded as the empty string across all 12{,}032 rows,
while the predicted answer field was populated. This is a capture gap
in the evaluation-generation pipeline, not model behaviour: the
responses were produced and the text was never persisted. Because the
trace accessor never returns an empty string
(Section~\ref{sec:arch-overview}), the judge was shown a 600--1200
character worked solution on one side and a synthesised one-line stub
on the other.

The defect was found while sizing the roster for a \emph{different}
experiment. It is reported here because it qualifies numbers already
written down, and because what survives it survives more strongly.

The roster splits perfectly on response format, and format tracks
capability: the four models with reasoning text score $57.7\%$ to
$78.4\%$ on the 241 judged Math questions, the four without score
$7.5\%$ to $53.1\%$, and the two groups do not overlap. Trace presence
and capability are therefore collinear by construction of this roster.

\textsc{measured.} On mixed pairs --- one traced response against one
stub --- the judge picks the side with text $863/868 = 99.4\%$ of the
time under \texttt{MAIN} and $874/880 = 99.3\%$ under
\texttt{GENERIC} (SE $0.3\%$), while ground truth says the traced side
is right on only $593/677 = 87.6\%$ of them. The excess preference is
$+11.8$ percentage points. \textbf{Mixed pairs are $868/1736 = 50\%$ of
all Run~A comparisons.} Half the pilot arena was decided by which side
had text at all. That is not a judge failure --- it is the only sane
response to an empty string --- but those verdicts carry no information
about model quality.

\paragraph{What this qualifies.}
The $88.8\%$ aggregate of Section~\ref{sec:res-correctness}: half its
support is format-decided, and the format preference agrees with the
key $87.6\%$ of the time \emph{by roster construction}. And the
aggregate rank correlations against ground-truth accuracy: ranking all
four traced models above all four untraced ones reproduces the
ground-truth top-four/bottom-four split for free, so those correlations
are substantially purchased by the artifact. Unaffected entirely are
the granularity and domain-screen results of
Section~\ref{sec:res-partition}, which use ground-truth accuracy only
and never touch a verdict.

\paragraph{What survives, and survives stronger.}
Within-stratum comparisons are unaffected, because both sides share a
format. Re-running the trace stratification with the capability gap
controlled strengthens the inversion of
Section~\ref{sec:res-trace-inversion} rather than dissolving it.

\begin{table}[h]
\centering
\caption{Answer-key agreement by trace stratum and ground-truth
accuracy gap (Run~A, \texttt{MAIN}). Cell counts in parentheses.}
\label{tab:capture-survives}
\begin{tabular}{lccc}
\toprule
Stratum & gap $<20$\,pp & gap 20--40\,pp & gap $\geq 40$\,pp \\
\midrule
trace $\times$ trace       & \textbf{83.2\%} (173) & 88.1\% (42) & --- \\
trace $\times$ no-trace    & 74.6\% (67)  & 83.3\% (186) & 94.6\% (407) \\
no-trace $\times$ no-trace & \textbf{97.0\%} (33) & 94.4\% (89) & 84.4\% (32) \\
\bottomrule
\end{tabular}
\end{table}

At matched capability gap ($<20$\,pp) bare-letter pairs reach $97.0\%$
against $83.2\%$ for reasoning pairs --- a $13.8$-point gap with
capability controlled. ``Reasoning text degrades the judge's agreement
with the key'' therefore survives as a within-format result rather than
a cross-format one. The worst cell in the table is the mixed stratum at
matched capability ($74.6\%$): near-equal models where the judge takes
the traced side almost always and is therefore wrong whenever the
untraced model happened to be right.

\paragraph{The repair.}
\textsc{measured.} The cause was in the ingest loader, and it was not
confined to four models: it dropped stored response text for 34 of the
corpus's 47 evaluated models. Re-ingest backfilled 356{,}402 rows and
took the corpus from 13 models with usable traces to 37. One of those
sits in a different question-identifier space and is excluded
(Section~\ref{sec:res-law}), so the maximum valid roster is 36 models.
Run~B is on the repaired corpus and on a length-matched band, so the
format artifact is controlled by construction there rather than
qualified after the fact.

\subsection{Rubric-Scoped and Rubric-Free Judges Reach the Same Verdict}\label{sec:res-generic-judge}

\textsc{measured.} Run~A's two conditions judged the same 1{,}736
(question, model-pair) comparisons, with complete overlap. Pairing them
row by row:

\begin{quote}
identical verdict: $1529/1736 = 88.1\%$ \\
disagreements: 207, of which 197 ($95\%$) involve a \texttt{TIE} on one
side \\
\textbf{actual winner flips: $10$ of $1736$, i.e.\ $99.4\%$ agreement
on winners}
\end{quote}

\textsc{measured}, and it rules out the obvious explanation: the rubric
\emph{is} read. The fraction of each rationale's vocabulary that also
appears in that cell's induced rubric is $0.200$ (IQR $[0.115, 0.280]$)
under \texttt{MAIN} against $0.138$ (IQR $[0.077, 0.200]$) under
\texttt{GENERIC\_JUDGE}, which never sees the rubric and therefore
measures topical coincidence between a Math rationale and a Math
rubric. \texttt{MAIN} sits $45\%$ above that baseline with the
interquartile ranges offset throughout. Mean rationale length is 54
words against 49.

\textsc{argued}, from those two measurements together: the verdict is
\textbf{rationale-conditioned, not rationale-driven}. For the
alternative --- that the judge applies the criteria and coincidentally
agrees --- to hold, a criteria-guided and a criteria-free judge would
have to reach identical winners on 1{,}726 of 1{,}736 comparisons. The
simpler reading is that the verdict is fixed before the rubric enters,
and the rubric shapes the justification afterwards.

Nor are the rubrics themselves generic:
Section~\ref{sec:res-rubric-null} shows they are measurably
cell-specific against a randomised null
($p = 1.21 \times 10^{-6}$). The specificity is real; on this corpus
it is inert at the verdict.

Two consequences. First, \textbf{a per-cell null between these two arms
is not evidence of parity} --- it is a comparison whose arms barely
differ, and no amount of statistical power repairs that. Second, and
this is a methodological result rather than an empirical one: the
precondition for any future A/B judging arm is to \textbf{measure
verdict agreement first}. It is one query against the verdict store,
and it decides whether the arm is worth ten thousand judge calls.

Caveat: lexical overlap demonstrates \emph{reproduction} of a cell's
vocabulary, not \emph{application} of its criteria. A verdict schema
carrying an explicit rule identifier would settle it directly --- a
cited rule cannot arise from topical coincidence --- and would also
count, for free, how many verdicts cite no rule at all and are
therefore decided on priors. That change alters the schema and so must
precede any further arena run.

\subsection{Tie Policy, Sampler, and the Stability of \texorpdfstring{$\rho$}{rho}}\label{sec:res-tie-policy}

\textsc{measured}, and it bounds every rank correlation in this thesis.
Tie handling moves Spearman $\rho$ by up to $\pm 0.029$ on this
pipeline, and the direction is \emph{not} consistent across conditions:
in Table~\ref{tab:c5-rho} of the next subsection, the per-leaf arm
gains $0.029$ from dropping ties while the generic arm loses $0.028$.

The sampler contributes a second source of movement. Two runs of the
same nominal condition --- the 12-model Math per-leaf arm --- span
$0.860$ to $0.916$ across the combinations of query sampler and tie
convention.

\textsc{argued}: absolute $\rho$ from this pipeline is not stable at
the precision a three-decimal figure implies, and should not be quoted
as though it were. Paired contrasts computed within a single run, on an
identical question set, under a stated tie convention, are stable ---
which is exactly what Run~B, next, was designed to produce. This is the
evidence behind rule~D5 of Section~\ref{sec:meth-discipline}, and it is
why every $\rho$ in this chapter appears with its tie convention
attached.

\subsection{The No-Partition Baseline}\label{sec:res-c5}

Removing the partition changes nothing the high-power metric can see;
the rank-correlation gap depends on tie scoring and, under the
registered rule, formally failed the null in the direction of the
partition-free arm. That composite is unpacked below.

Section~\ref{sec:res-generic-judge} varies the rubric while holding the
partition fixed. It cannot test the partition, because both of its arms
replay the same leaf-scoped triples. Run~B varies the partition itself:
one arm judges each question under its own leaf's rubric and rolls the
per-leaf fits up by pooling sufficient statistics; the other judges the
same questions under a generic MT-Bench-style pairwise prompt and fits
Bradley--Terry directly at the domain. Unlike the original MT-Bench
protocol, both arms judge an identical question set, so the comparison
is paired and isolates the partition.

The reading was fixed in writing before the comparison existed. The
primary outcome is Spearman $\rho$ against ground-truth accuracy per
domain; the secondary is per-verdict answer-key agreement, which has
far more power but measures the judge rather than the grouping. The
decisive difference on Math was fixed at two Spearman grid steps: at
$M = 12$ the grid is $6/(12 \cdot 143) = 0.0035$, so
$|\Delta\rho| \geq 0.007$ counts as a real difference and anything
below it is read as flat. Math was designated the \emph{null} arm in
advance, on the ground that the band's models order identically in Math
and globally (Section~\ref{sec:res-law}), so a per-leaf advantage there
would indicate an artifact rather than a finding. Four void conditions
were named: the arms must judge the same question set, the confidence
gate must be disabled in both, both must use the same query sampler,
and ties must be defined identically across the two verdict formats.

\textsc{measured}, secondary outcome: the two arms are
indistinguishable. Answer-key agreement is $94.2\%$ (SE $0.7\%$) for
the per-leaf arm and $94.7\%$ (SE $0.6\%$) for the generic arm.

\textsc{measured}, primary outcome, and it depends on the tie
convention:

\begin{table}[h]
\centering
\caption{Run~B, Math. Spearman $\rho$ between fitted Bradley--Terry
$\theta$ and ground-truth accuracy, under both tie conventions. The
pre-registered decisive difference is $|\Delta\rho| \geq 0.007$.}
\label{tab:c5-rho}
\begin{tabular}{lrrr}
\toprule
Tie convention & per-leaf (\texttt{MAIN}) & generic (\texttt{C5}) & $\Delta\rho$ \\
\midrule
ties half-weighted & 0.887 & 0.951 & $+0.064$ \\
ties dropped       & 0.916 & 0.923 & $+0.007$ \\
\bottomrule
\end{tabular}
\end{table}

Half-weighting ties puts the generic arm ahead by nine times the
decisive threshold. Dropping tied comparisons collapses the gap to
exactly the threshold. Nothing about the verdicts changes between the
two rows --- only how a tie is scored.

The quoted values use an unpenalised pooled MM fit
(\texttt{tools/analysis/tie\_policy\_sweep.py}); the run's own export,
which applies a Jeffreys prior, gives $0.874$ and $0.965$ under
half-weighting --- a third convention, and it is named rather than
left implicit.\footnote{Named because an unstated analysis convention
is how figures become irreproducible; the corrections register
(Appendix~\ref{app:corrections-register}) records the instances this
project has already paid for.}

\textsc{measured}, against the registered rule: the pre-registration
fixed $|\Delta\rho| \ge 0.007$ as a \emph{failed} null, and both
conventions meet it --- $+0.064$ half-weighted, $+0.007$ with ties
dropped, both in favour of the partition-free arm. Under the rule as
written, \textbf{the Math null failed}.

\textsc{argued}, and after the fact: the failure is a tie-scoring
artifact rather than a partition effect. Nothing about the verdicts
changes between the two rows; the high-power agreement statistic, which
resolves differences of a few points on thousands of verdicts, sees
nothing; and the surviving gap sits exactly on the registered boundary
rather than beyond it. But the convention that shrinks the gap to the
boundary --- rule D5, both tie policies --- was adopted only after
these data existed, so it cannot re-read the registered verdict; it can
only motivate a confirmatory design in which the tie convention is
fixed in advance. Prediction~1 is therefore recorded as \emph{failed
under the registered rule, with a post-hoc attribution to tie scoring},
not as held. At the roll-up level both arms are fitted on the identical
250-question set, so the registered two-sided disattenuation is the
same factor on both sides and cancels in this contrast; it is therefore
not shown.

What Run~B does \emph{not} test is worth stating with the same
precision, because pooling makes it true by construction: it does not
test whether the partition helps aggregation, since both arms are
identical in estimation. It tests rubric conditioning at matched
everything-else. And it does not test the partition where the partition
has something to be right about; Math is the null arm by design, and
the corresponding test is Run~C.


\paragraph{Run~C: the law result.} The paired design held: both arms
judged the same 246 questions (\textsc{main} 2{,}499 comparisons, ties
19.8\%; \textsc{c5} 2{,}479, ties 17.8\%). \textsc{measured}, on
answer-key agreement: \textsc{main} 982/1{,}116 = 88.0\% against
\textsc{c5} 1{,}039/1{,}206 = 86.2\% (SE 1.0\% each). \textsc{measured},
on the pooled rank correlation against ground truth, both tie
conventions: half-weighted, \textsc{main} 0.9510 against \textsc{c5}
0.9510 ($\Delta = 0.0000$); ties dropped, \textsc{main} 0.9650 against
\textsc{c5} 0.9441 ($\Delta = +0.0210$).

Run~C used the twelve-model roster, so the grid step is $0.0035$ and
the two-step decisive threshold is $|\Delta\rho| \geq 0.007$. The
registered prediction for law --- per-leaf judging beats partition-free
judging, if the partition helps at all --- is met in direction under
both conventions and in magnitude under one: ties dropped,
\textsc{main} leads by three times the threshold; half-weighted, the
arms are exactly tied. Unlike the Math contrast, the gap never changes
sides: the law result is parity or a per-leaf advantage, not a
reversal. And unlike Math, the both-conventions rule (D5) predates this
run, so dual reporting here is pre-specified rather than post hoc. The
answer-key-agreement gap ($+1.8$ points, about $1.3$ combined standard
errors) leans the same way without being independently decisive. The
prediction pair therefore closes as registered: flat on Math, where the
roster's models do not reorder, and a convention-sensitive but
never-reversed per-leaf advantage on law, where they do.

\paragraph{The near-clone check, evaluable for the first time.} The
bootstrap floor guaranteed the two Claude~3.5 Sonnet harnesses direct
comparisons --- 24 per arm, where Run~B gave the pair zero. Of those,
11 of 24 were ties in each arm, a 46\% pairwise tie rate against the
roster average of roughly 19\%: the judge declines to separate two
harnesses of nominally the same model far more often than it declines
elsewhere --- the resolution behaviour the pre-registered check
demanded (prediction P2, met).

\subsection{Identification Cost}\label{sec:res-cost}

\textsc{measured.} On Run~A, per-cell identification --- a connected
comparison graph over the roster --- is reached at \textbf{7 model
pairs of 28}, which is a spanning tree and is what the theory predicts,
at between $32.9$ and $45$ comparisons per cell. The cost \emph{varies}
across cells, which clears the void condition the pre-registration
named in advance: an earlier run's uniform figure was a scheduler
allocation rather than a measured cost, and a second constant would
have meant the scheduler was still capping.

The minimum, $32.9$, sits inside the $\leq 40$ band fixed in advance,
so leaf-level judging is affordable: roughly 2{,}900 to 3{,}900 judge
calls for 87 cells against a budget of about 10{,}000.

Caveat: measured at 8 models, hence 28 pairs. At 11 models there are 55
pairs and the arithmetic was not redone.

\subsection{What Section~\ref{sec:res-judge} Establishes}\label{sec:res-judge-summary}

\textsc{measured}, composite. The judge \textbf{decides on answer
correctness} ($88.8\%$ agreement with the key on discriminable items,
with half that support format-decided); \textbf{reads} the
leaf-specific rubric (rationale--rubric overlap $0.200$ against a
$0.138$ baseline that never sees it); \textbf{reproduces its
vocabulary} in the justification; and \textbf{reaches the same verdict
as a rubric-free judge} ($99.4\%$ identical winners) and as a
partition-free one ($94.2\%$ against $94.7\%$ answer-key agreement).
When correctness cannot discriminate, the tie rate triples. Each
component has a control sharing every confound except the treatment.

\textsc{argued}, from all of the above: \textbf{on a multiple-choice
corpus with a verifiable key, a capable judge can shortcut to
correctness and bypass the partition entirely.}

This is a hypothesis with strong support, not a demonstration. It is an
inference from Math pilots, one of which carries a named format
artifact, and the decisive experiment --- withholding the
multiple-choice options from the judge while holding the question set
fixed --- has not been run. Its predictions are fixed in writing
(Appendix~\ref{app:tuning-protocol}) and it is blocked on a
trace-complete roster: without options \emph{and} without stored
reasoning, an untraced-versus-untraced comparison is uninterpretable
rather than a chance-level null.

\section{What the Partition Does}\label{sec:res-partition}

\subsection{Link 1: The Construction Meets Its Requirements}\label{sec:res-construction}

\textsc{measured.} The frozen artifact has 154 nodes, 87 leaves,
maximum depth 6, $J = 0.253129$, and conserved query mass $8299.00$
with no mass warnings. All 87 leaves are labelled and carry a judge
rubric. Leaf populations have median 88 queries, IQR $[71, 118]$,
minimum 54 and maximum 396, with one leaf ending below the birth floor
of 55 for the reason given in Section~\ref{sec:arch-limits}.

\textsc{measured.} The construction reaches a certified fixed point at
iteration 10: structural edits $0$, maximum
$1 - \cos(\hat{\boldsymbol{\mu}})$ across all nodes
$1.11 \times 10^{-16}$, maximum relative change in concentration
$4.25 \times 10^{-13}$, against a tolerance of $10^{-6}$; iterations 8,
9 and 10 are identical on every recorded column. One further
application of the operator would change nothing the artifact consists
of --- not the structure, not the partition, and not the fitted node
parameters, which matters because held-out queries route through those
parameters.

\begin{figure}[h]
\centering
\fbox{\parbox{0.9\textwidth}{\centering\vspace{0.6cm}
Figure --- Leaf-size distribution (placeholder)\\
\footnotesize Histogram of leaf query counts in the frozen snapshot,
vertical reference line at the birth floor $n_{\min} = 55$, median and
IQR marked; data: final iteration of the construction snapshot
export.\vspace{0.6cm}}}
\caption{Distribution of leaf query counts in the frozen snapshot, with
the birth floor $n_{\min} = 55$ marked. The floor is enforced at split
time on the routed partition and is not an invariant of the final
tree; one leaf of 87 sits below it.}
\label{fig:leaf-size-hist}
\end{figure}

\paragraph{The acceptance gate is stability-neutral and $J$-neutral,
and both were measured rather than assumed.}
\textsc{measured}, and both results are negative. The gate's effect on
run-to-run structure is an additive offset of about $4.4$ leaves with
no detectable change in dispersion; the statistical mechanics of that
reading are the second instance of rule~D3
(Section~\ref{sec:meth-discipline}). And the gate does not raise the
objective it gates: $J$ is higher on 2 of 5 seeds, sign test
$p = 0.812$, median $\Delta J = -0.0001$; the apparent gain was seed 42
alone.

\textsc{argued}: the gate is adopted because a fixed absolute
tolerance is incoherent against a $13.1\times$ spread in the standard
error of the quantity it thresholds, and it is adopted \emph{after}
both of its natural empirical supports were tested and refused --- the
transferable statement of which is Section~\ref{sec:arch-dag}. A
corroboration that was tested and failed is stronger evidence about a
rule's justification than one that was never tested at all.

Caveat: the five-seed sweep's leaf counts predate a correction to the
splitter's routing check and are \textsc{void as to values} (marker
defined in Section~\ref{sec:arch-dag}). The analysis design and its
two negative results are unaffected --- they are comparisons internal
to the sweep --- but the underlying counts need re-deriving.

\subsection{Link 2: Rubric Specificity Against a Randomised Null}\label{sec:res-rubric-null}

The question is whether coherent cells produce more specific rubrics
than incoherent ones. It is offline: it needs no arena. It was
pre-registered with four directional predictions, the confound
direction named per measure, and decision bands fixed in advance.

\textsc{measured.} Rubrics were induced by the identical procedure on
the 87 leaves and on 13 size-matched random cells drawn from the same
corpus. The specificity measure --- the fraction of a rubric's
vocabulary found in its own cell's queries, above what it shares with
other cells --- separates the two arms decisively: Mann--Whitney $U$,
one-sided, $p = 1.21 \times 10^{-6}$, with medians $0.138$ against
$0.012$. The result is robust across all three definitions of the
per-cell statistic that were considered
($1.2\times10^{-6}$, $2.1\times10^{-6}$, $7.6\times10^{-6}$).

\textsc{measured.} All four pre-registered discriminators fire in the
predicted directions.

\begin{table}[h]
\centering
\caption{Rubric-specificity discriminators, leaf cells against
size-matched random cells. The predicted direction was fixed in advance
for each row.}
\label{tab:rubric-null}
\begin{tabular}{lrrl}
\toprule
Measure & 87 leaf cells & 13 random cells & predicted \\
\midrule
own-cell specificity (median) & \textbf{0.138} & 0.012 & leaf higher \\
pairwise Jaccard between rubrics & 0.0754 & 0.1187 & leaf lower \\
terms present in $\geq 90\%$ of rubrics & \textbf{0} & 6 & leaf $\approx 0$, random $>0$ \\
cells with positive specificity & 81/87 (93\%) & 7/13 & leaf higher \\
\bottomrule
\end{tabular}
\end{table}

The sharpest of these is the third. A random cell contains chemistry,
law and history at once, so its rubric must fall back on domain-general
language; a generic core is present across the random rubrics and
absent from every leaf rubric. Had both been near zero the instrument
could not see the difference and the test would have been inconclusive
rather than negative --- a distinction fixed in advance.

The random arm's absolute specificity of $0.012$ falls inside the
$\leq 0.05$ band fixed before the run as meaning ``coherence is the
cause''. An arm-size control was run because 13 against 87 is
unbalanced: 13 leaf rubrics subsampled give $0.134$, and a 200-draw
bootstrap over subsamples gives $0.138$, so the treatment arm's value
is not a function of its size.

Three caveats travel with this result and none of them is optional.
First, the measure is \textbf{lexical throughout}: it demonstrates
reproduction of a cell's vocabulary, not application of its criteria.
Second, the treatment bundles coherent membership with an informative
topical label, and the null cells were deliberately labelled
\texttt{Cluster NN}; the claim the design supports is therefore
``a coherent cell together with its label causes specificity'', not
coherence alone. Third, 13 random cells is a small null arm.

One consequence must be stated here rather than left to
Chapter~\ref{ch:discussion}: these are the same rubrics that
Section~\ref{sec:res-generic-judge} shows leave the judge's winner
unchanged on $99.4\%$ of comparisons. Link~2 holds --- the cells
produce measurably specific rubrics --- and on this corpus that
specificity is inert at the verdict.

\todo[inline]{Two items before this section ships. (i) The supporting
artifacts for this result --- the null-arm generation script and its
test harness --- are untracked in version control; a result whose
artifacts are untracked is one clean checkout from being unsupported.
(ii) Two pre-registration documents state different decision bands
($\leq 0.05$ / $\geq 0.10$ in one, $\geq 0.05$ with $\geq 75\%$ of
cells positive in the other) and neither says whether they govern the
same quantity. Resolve which band applies to which statistic before
either is quoted in the final text.}

\subsection{Link 3's Premise Fails, and It Was Pre-Registered}\label{sec:res-granularity}

\textsc{measured}: cells do not carry different true model rankings on
this corpus, at any granularity from 14 to 152 cells. The analysis was
fixed in writing before it was run --- three outcomes named in
advance, ``observed $\rho$ is primary'' fixed in advance --- and that
registration is the only reason the result is readable, because the
disattenuated column moves the opposite way and clips past $1.0$.

The question is whether cells carry different \emph{true} model
rankings: for every pair of cells, the Spearman correlation between the
two cells' ground-truth accuracy rankings of the roster. If cells
differ in what they measure, that correlation should fall as cells get
finer.

\begin{table}[h]
\centering
\caption{Observed between-cell Spearman $\rho$ of ground-truth model
rankings, across an eleven-fold change in granularity. The 88-leaf and
87-leaf rows differ only in the router used to assign queries and are
an explicit control.}
\label{tab:granularity}
\begin{tabular}{lrrr}
\toprule
Partition & cells & cell pairs & median $\rho$ \\
\midrule
MMLU-Pro editorial domains & 14 & 91 & 0.929 \\
induced leaves (biased router) & 88 & 3828 & 0.922 \\
induced leaves (corrected router) & 87 & 3741 & 0.922 \\
induced leaves, finer setting & 152 & 11476 & 0.905 \\
\bottomrule
\end{tabular}
\par\smallskip
\footnotesize One reproduction note, stated because the convention
matters: the leaf-level $0.922$ reproduces on the snapshot's own
construction assignments over the full corpus; recomputing it on
re-routed, reserved-only assignments gives $0.884$. Both are
defensible readings of ``the ranking within a leaf''; they are not the
same quantity, and the value quoted here is the first. The conclusion
--- flatness across granularity --- is unchanged under either, since
the comparison is between rows computed the same way.
\end{table}

It does not fall. Interquartile ranges overlap throughout. The 87- and
88-leaf rows are identical to three decimals although the routing
correction between them is worth $+10.3\%$ on the construction
objective at matched granularity --- so discriminative power is flat
across \emph{routers} as well as across granularities. Centring out
each model's global mean leaves residual correlations of $-0.119$,
$-0.024$, $-0.024$ and $-0.024$, all within $1.5$ to $3.6$ times the
mechanical centring null $-1/(C-1)$; the correct statement is
``indistinguishable from the centring null'', not ``zero''.

The join underlying this analysis was verified before it was computed,
per rule~D2: exact, $11769/11769$; all roster models present in every
cell of every tree; no cells dropped.

\textsc{measured}, at the corrected roster. The flatness is not an
artifact of the 8-model pilot roster. Re-derived at the corrected band
it holds: $0.929$ and $0.922$ at 8 models, and $0.973$ and $0.936$ at
the 11-model band.

\paragraph{Three bounds, all pre-registered.}
\emph{Power, not proof.} Per-model per-cell accuracy has a standard
error of about $2.1$ percentage points at 14 domains, $5.8$ at 87
leaves and $6.9$ at 152. The null is consistent with ``no cell-specific
ranking signal'' and with ``a signal below about 6 points at leaf
granularity''; this analysis cannot separate them. The 14-domain row is
the stronger evidence precisely because it is the low-noise condition
and its residual is also at null.
\emph{Roster.} The pilot roster spans $54.1$ accuracy points, so
$\rho \approx 0.92$ was largely measuring that the strongest model beats
the weakest in every cell. A roster clustered within about 10 accuracy
points is untested and is the version of this question that could still
come out positive.
\emph{Ceiling, not prediction.} This is measured on each cell's own
constituent queries and on ground-truth accuracy --- never on routed
held-out queries and never on judge verdicts. It \emph{bounds} what an
arena could discover in these cells; an arena can show less
discrimination than this, never more.

\textsc{argued}: this is a property of the corpus and the roster, not
of the method. MMLU-Pro is saturated for a roster of this spread, and
cell identity contributes essentially nothing to model ranking beyond
global capability ordering, at any granularity from 14 to 152, under
either router. It says nothing about judge quality, which is what
Section~\ref{sec:res-judge} measures, and a reader must not be allowed
to collapse the two.

\begin{figure}[h]
\centering
\fbox{\parbox{0.9\textwidth}{\centering\vspace{0.6cm}
Figure --- Between-cell rank agreement across granularity (placeholder)\\
\footnotesize Four box plots (14 / 88 / 87 / 152 cells) of the
between-cell Spearman distribution, with the mechanical centring null
drawn as a reference line on a companion panel of centred residuals;
data: the pre-registered granularity analysis export.\vspace{0.6cm}}}
\caption{Between-cell agreement of ground-truth model rankings across
an eleven-fold change in granularity, with the centred-residual panel
alongside. The overlapping interquartile ranges, not the medians, carry
the result.}
\label{fig:granularity-flat}
\end{figure}

\subsection{The Domain Screen: Law Reorders, Math Does Not}\label{sec:res-law}

An arena can only reveal cell-specific capability where cell-specific
capability exists. The screen is offline: for each MMLU-Pro domain,
correlate the roster's ground-truth accuracy ranking within that domain
against its global ranking, and compare against a size-matched
permutation null over random subsets of the same reserved pool.

\textsc{measured}, at the 11-model band, 1{,}500-draw null, seed 42,
under two roster policies: \textbf{law} reorders
($n = 287$, $\rho = 0.955$, $p = 0.008$) and \textbf{Math} does not
($n = 393$, $\rho = 1.000$, $p = 1.000$). At this band the models order
identically in Math and globally, so an arena there has nothing to
detect --- which is why Math is the pre-registered null arm of Run~B
and law is the domain of Run~C.

\textsc{measured}. The process history is itself instructive:
\textbf{this conclusion reversed twice before settling, and both
reversals had the same cause.}

At 8 models, Math was statistically indistinguishable from a same-size
random subset ($p = 0.217$) and only law, philosophy and history cleared
the null; the recommendation was law. After the trace backfill the
screen was re-run at 37 models, Math became significant
($p = 0.001$) and seven domains cleared the null; the recommendation
flipped to Math. That flip rested on a model-intersection poisoned by a
single model sitting in a different question-identifier space, whose
inclusion collapsed law from 287 questions to 20 and made law look
unaffordable. Excluding it restores law to its full 287 and yields the
11-model band, at which Math scores $\rho = 1.000$ and the original
recommendation returns.

\textsc{argued}: both inversions came from evaluating a ranking
statistic at a model count that did not match the decision it was
informing. The screen must be computed at the roster the arena will
actually run. That rule is worth more than either recommendation, and
it is reported here rather than quietly applied.

Second caveat, and it bounds the whole section: both screens use
MMLU-Pro's native category labels, not induced cells. They \emph{bound}
what a partition of a domain could show --- a domain whose own label
carries no reordering is unlikely to contain cells that do --- but they
do not show that induced cells inside law separate. That remains the
test, and it is Run~C.

\subsection{Polyhierarchy: A Measured Negative Result}\label{subsec:polyhierarchy}\label{sec:meth-topology}\label{sec:poly-negative}

The construction produces a strict tree: every node below the root has
exactly one parent.
An earlier version of the pipeline admitted polyhierarchy through a
dedicated operator. This section records that operator, the evidence
that retired it, and the invariant that replaced it.

Two claims have to be kept apart.
Construction-time routing is weighted multi-membership
(Appendix~\ref{app:dag-formalism}), which lets a \emph{query} belong to
several leaves.
Polyhierarchy is a different claim: it lets a \emph{concept} belong to
several branches, through node-level topology edits.
The first is retained; the second is not.

\textbf{The operator that was tested.}
A cross-link was a growth edit: an existing node $n$ gained a second
parent $h$ when $h$'s residual pool demonstrably fitted $n$.
The motivating case is a genuinely cross-domain query---a
biostatistics question between Mathematics and Biology---which under a
strict tree can only pick one branch, fail to specialise there, and
residualise at the domain anchor.
Candidates were ranked by the fraction of $h$'s residuals they
\emph{captured}, with capture applying the runtime descent gate
itself, so that ``captured'' meant ``will actually descend at $h$ once
the edge exists,'' not proxy similarity. The capture criterion, the
guard list, and the acceptance threshold are specified in
Appendix~\ref{app:cross-link-operator}.

\textbf{What the measurements showed.}
\textsc{measured.} Across 42 cross-link proposals on the canonical
corpus, the capture fraction $f$ that the operator uses to rank
candidates was uncorrelated with the objective it must satisfy:
$r(f, \Delta J) = +0.085$, with complete distributional overlap between
accepted and rejected proposals (median $f$ of $0.802$ accepted against
$0.818$ rejected).
The effect sizes were an order of magnitude apart: median $\Delta J$
was $6.9 \times 10^{-5}$ for cross-links against $5.3 \times 10^{-4}$
for splits.
Admission was therefore decided by evaluation order rather than by
evidence, which the proposal sequence makes explicit: one node was
accepted at four hosts and rejected at two, and acquired five parents
by iteration~1.

\textbf{Why a partition objective cannot decide this.}
$J$ scores a partition of query mass into cells.
Giving a node a second parent changes which paths reach its cell, but
it leaves the cell contents almost unchanged, so $\Delta J$ carries
little signal about whether the edge is semantically right.
A partition-based objective is structurally indifferent to
polyhierarchy, and no threshold on such an objective can separate good
cross-links from bad ones.

\textbf{The decision, and what it costs.}
The operator was removed rather than retuned.
The multi-membership claim survives that removal: the multi-leaf rate
of held-out queries was $0.1549$ with cross-links and $0.1518$ without,
so queries that legitimately span domains still reach several leaves
through the routing beam and the membership floor.
What is given up is the \emph{concept}-level claim, which the evidence
above could not support in the first place.

\textbf{The single-parent invariant.}
Removing the generator is not by itself sufficient, because fusion also
redirects parent edges: fusing two nodes that sit under different
parents leaves the survivor holding both, which reintroduces
multi-parent nodes with no cross-link operator present.
Fusion proposals are therefore restricted to pairs sharing a parent,
which makes every topology operator a map from a tree to a tree.
The structural-integrity gates of Appendix~\ref{app:tuning-protocol}
check the invariant directly by counting multi-parent nodes, which must
be zero.

Evaluation-time multi-leaf assignment of held-out queries is a separate
mechanism again. It applies the same routing read-only to the frozen
snapshot, capped at \texttt{maxLeafAssignments}, and so does not modify
the topology; what it determines is which leaves a given query
contributes comparisons to during the arena run
(Section~\ref{sec:arch-arena}).

\section{Chapter Summary}\label{sec:res-summary}

Stated against the criteria fixed before the data, with no
interpretation.

\textbf{RQ1, what decides the verdicts.} On a corpus with a verifiable
key, the judge's verdict tracks answer correctness. Agreement with the
key is $88.8\%$ on discriminable items; the tie rate triples where the
key stops discriminating; agreement is highest where neither response
carries reasoning and lowest where both do, and that inversion
strengthens to $97.0\%$ against $83.2\%$ once capability is controlled;
and half of the pilot run's comparisons were decided by a response-text
capture gap, which is reported as a result with what it qualifies named.

\textbf{RQ1, rank recovery.} Both judges recover the ground-truth
ordering to within about one adjacent transposition, but absolute rank
correlations from this pipeline move by up to $\pm0.029$ under tie
convention and span $0.860$--$0.916$ across sampler and tie-convention
choices for a single condition. Only paired contrasts are quoted.

\textbf{RQ2, the partition.} Under the registered rule the Math null
\emph{failed} in the direction of the partition-free arm; the
attribution to tie scoring is post hoc. On
an identical 250-question set, per-leaf and generic partition-free
judging agree to $94.2\%$ against $94.7\%$ on answer-key agreement; the
Spearman contrast is $+0.064$ under half-weighted ties and $+0.007$ ---
exactly the pre-registered decisive threshold --- under dropped ties.
Math was designated the null arm in advance. The corresponding test on
law is Run~C and is outstanding.

\textbf{Link 1.} Met. Certified fixed point at iteration 10; 87 leaves,
median size 88; the acceptance gate is stability-neutral and
$J$-neutral, both measured.

\textbf{Link 2.} Positive. Rubrics induced on coherent cells are
specific to those cells against a size-matched randomised null,
$p = 1.21 \times 10^{-6}$, with all four pre-registered discriminators
firing in the predicted directions. Lexical measure only.

\textbf{Link 3.} Premise refused, under pre-registration, before the
arena ran. Between-cell agreement of true rankings is flat across an
eleven-fold change in granularity and across two routers, bounded by
power, by roster spread, and by the fact that it is a ceiling rather
than a prediction.

\textbf{Polyhierarchy.} Operator built, measured, retired. The capture
score it ranked candidates by is uncorrelated with the objective it
must satisfy, and a partition objective is structurally indifferent to
polyhierarchy. Query-level multi-membership survives; the
concept-level claim does not.
